% STATUS: CANONICAL SPECIFICATION; no theorem is asserted or proved here.
\section{Canonical model contract}
\label{sec:canonical}
The explicit initialization prompt is authoritative. This fragment specifies
the main models; formulas in the archived blueprint remain research candidates.

\subsection{Main demand-modulated model}
For $i\in\Zn$, define the vector field $F^{\DM}$ by
\begin{equation}
\label{eq:canonical-dm}
\dot{x}_i=a x_i-kx_i^2+\frac{c x_{i-1}}{1+x_i}
             -\frac{g x_i}{1+x_{i+q}}-h x_i x_{i-q}.
\tag{$\DM_{n,q}$}
\end{equation}
The variables are dimensionless normalized activities. The preferred term is
\emph{demand-modulated model}. The positive input is called demand-modulated
predecessor input; no standard Holling type-II classification is adopted.
The intended initial conditions are $x(0)\in\Orth$. For local statements and
boundary bifurcations, use the specified open extension domain
$\Domain=(-1,\infty)^n$. This domain choice asserts no invariance theorem.

\subsection{Primary matched bilinear comparator}
Define $F^{\MB}$ by
\begin{equation}
\label{eq:canonical-mb}
\dot{y}_i=a y_i-k y_i^2+c y_i y_{i-1}
                 -g y_i y_{i+q}-h y_i y_{i-q},\qquad i\in\Zn.
\tag{$\MB_{n,q}$}
\end{equation}
This is the topology- and parameter-count-matched comparator in the main
program. Its intended initial data are nonnegative; the polynomial vector
field is specified on $\R^n$. No existence, invariance, boundedness or
stability conclusion is included in this contract. Matching parameter names
does not settle a quantitative calibration between the two laws.

\subsection{Parameter and indexing contract}
\[
n\in\Z,\quad n\geq5,\qquad q\in\{1,\ldots,n-1\},\qquad
\param=(a,k,c,g,h)\in(0,\infty)^5.
\]
All indices are modulo $n$, with representatives $0,\ldots,n-1$.
Arbitrary $q$ is allowed. The interpretation as one connected step-$q$ cycle
requires $\gcd(n,q)=1$; this condition is not imposed on the general analytical
program. In particular, even $n$ is not excluded when $q=2$. Coincident offsets
do not remove terms from either defining equation.
The intended exact wave-number-selection program focuses on $q=2$.
The five-node Wu-Xing-inspired specialization is $n=5,q=2$.

\subsection{Archived secondary six-term comparator}
The following is a transcription of the blueprint's proposed arbitrary-$n,q$
extension, not a verified attribution to the original Lin--Han publication:
\begin{equation}
\label{eq:archived-lb}
\dot{y}_i=a y_i-b y_i^2+c y_i y_{i-1}-d y_i y_{i-q}
                    -g_{\mathrm c}y_i y_{i+1}-h y_i y_{i+q}.
\tag{$\LB_{n,q}$ archive}
\end{equation}
Its archived five-node equation is obtained by specifying $n=5,q=2$ in this
transcription. The six coefficient names are archival and are not a parameter
identification with $\DM$ or $\MB$. The original source, coefficient domain,
index orientation, and attribution remain to be checked before analysis or
use as a literature application.

\subsection{Scope boundary}
The main program covers homogeneous $\DM$, primary comparator $\MB$, their
synchronous equilibria and spectra, the planned bifurcations, and the planned
small-heterogeneity response theory. All requested P01--P18 tasks remain in
the project; final main-text versus supplement placement is open.
The heterogeneous coefficient assignment is not frozen by this contract.
The six-term model is archived for a possible secondary literature application,
never a replacement for $\MB$. Source-saturating replacements, local motif
theory, pathology extensions, adaptive set points, interventions, noise,
delays, data calibration, and prior Lyapunov-spectrum reanalysis are postponed
to other work. The master remains unchanged.
